\originalchapter{振动与物理摆}\label{ch:osc}
\originalsection{简谐运动与能量}
弹簧偏离平衡产生与位移成正比的恢复力, $m\ddot x=-kx$. 写 $\omega_0=\sqrt{k/m}$, 满足初值 $x(0)=x_0,\dot x(0)=v_0$ 的解为
\begin{equation}
x=x_0\cos\omega_0t+(v_0/\omega_0)\sin\omega_0t
=A\cos(\omega_0t+\phi).
\end{equation}
振幅 $A=\sqrt{x_0^2+(v_0/\omega_0)^2}$, 相位 $\phi$ 须按初值确定象限. 周期 $T=2\pi/\omega_0$, 频率 $f=1/T$, 角频率 $\omega_0=2\pi f$, 三者不能混用. 动能与弹性势能交替变化, 总能量 $E=kA^2/2$. 竖直弹簧的重力改变平衡位置, 在新平衡附近同样有 $\omega_0=\sqrt{k/m}$.

一般光滑稳定势能极小点附近, 第\ref{ch:energy}章给 $k_{\rm eff}=U''(x_e)$, 因而小振动频率 $\sqrt{k_{\rm eff}/m}$. 线性近似的条件是高阶势能项相对二次项很小. 不是所有有周期的运动都是简谐运动.

\originalsection{单摆与物理摆}
单摆长 $l$、质点质量 $m$, 关于悬点的惯量 $ml^2$, 重力矩 $-mgl\sin\theta$, 得 $\ddot\theta+(g/l)\sin\theta=0$. 小角 $|\theta|\ll1$ 时 $\sin\theta\simeq\theta$, 周期 $2\pi\sqrt{l/g}$. 物理摆为任意刚体绕固定水平轴摆动, 质心离轴 $d$, 惯量 $I_O$, 则 $I_O\ddot\theta=-Mgd\sin\theta$, 小角周期
\begin{equation}
T=2\pi\sqrt{I_O/(Mgd)}.
\end{equation}
均匀杆端悬 $I_O=ML^2/3$, $d=L/2$, 因而 $T=2\pi\sqrt{2L/(3g)}$. 质心恰在轴上 $d=0$ 时没有重力恢复矩, 上式不能当作普通摆周期.

有限振幅 $\theta_0$ 时从能量 $\dot\theta^2=(2g/l)(\cos\theta-\cos\theta_0)$ 得
\begin{equation}
T=4\sqrt{l/g}\int_0^{\pi/2}
\frac{\dd\psi}{\sqrt{1-k^2\sin^2\psi}},\qquad k=\sin(\theta_0/2).
\end{equation}
使用 $\sin(\theta/2)=k\sin\psi$ 消除了转向点奇点. 展开 $(1-k^2\sin^2\psi)^{-1/2}$ 并积分, 得 $T=2\pi\sqrt{l/g}[1+\theta_0^2/16+O(\theta_0^4)]$. 有限振幅周期较大, 接近上端不稳定平衡 $\theta_0\to\pi$ 时周期发散.

\originalsection{阻尼、驱动与共振}
线性阻尼振子 $m\ddot x+b\dot x+kx=0$. 写 $\gamma=b/(2m)$, 特征根 $s=-\gamma\pm\sqrt{\gamma^2-\omega_0^2}$. 欠阻尼 $\gamma<\omega_0$ 时 $x=e^{-\gamma t}(A\cos\omega_dt+B\sin\omega_dt)$, $\omega_d=\sqrt{\omega_0^2-\gamma^2}$. 临界阻尼为 $(A+Bt)e^{-\omega_0t}$; 过阻尼为两个负实指数. 阻尼能量变化为 $\dot E=-b\dot x^2\le0$, 欠阻尼的周期平均包络按 $e^{-2\gamma t}$ 衰减, 瞬时机械能不是单一严格指数.

有驱动力 $F_0\cos\omega t$ 时, 稳态解 $x=X\cos(\omega t-\delta)$, 代入方程并分别比较正弦余弦系数得
\begin{equation}
X=\frac{F_0}{\sqrt{(k-m\omega^2)^2+b^2\omega^2}},\qquad
\delta=\operatorname{atan2}(b\omega,k-m\omega^2).
\end{equation}
瞬态会按阻尼消失. 平均储能 $\langle E\rangle=(m\omega^2+k)X^2/4$, 平均输入与耗散功率均为 $b\omega^2X^2/2$. 位移振幅峰值在 $\omega^2=\omega_0^2-2\gamma^2$ (右侧为正时), 平均吸收功率峰值在 $\omega=\omega_0$. “共振频率”应说明观察的是哪一个量. 弱阻尼品质因子 $Q=m\omega_0/b=\omega_0/(2\gamma)$, 描述储能相对每弧度耗散的比值; $Q$ 大则共振峰窄.
\originalsection{方程求解与复合摆}
简谐方程的两个解 $\cos\omega_0t,\sin\omega_0t$ 线性独立, 任意初值可由它们组合满足. 两个具有相同初值的解之差 $z$ 初值为零; 其能量 $(m\dot z^2+kz^2)/2$ 恒为零, 所以 $z=0$. 这证明给出的解是完整初值解, 不只是一组试探解.

复数 $i^2=-1$, Euler 公式 $e^{i\phi}=\cos\phi+i\sin\phi$ 可把两实振荡合成一个指数. 常系数线性方程代 $e^{st}$ 得特征多项式; 共轭根 $s=-\gamma\pm i\omega_d$ 的实线性组合是前述指数乘正余弦. 重根时令 $x=e^{-\omega_0t}y$, 方程化为 $\ddot y=0$, 因而 $y=A+Bt$, 解释了临界阻尼的 $t$ 项.

驱动振动限于 $b>0$, 可令 $x=\Re(Ze^{i\omega t})$, 求得 $Z=F_0/(k-m\omega^2+ib\omega)$. 其模给振幅, 辐角给相位滞后, 与正文公式相同. 等价的实系数法令 $x=A\cos\omega t+B\sin\omega t$, 得 $(k-m\omega^2)A+b\omega B=F_0$, $(k-m\omega^2)B-b\omega A=0$, 解得 $A=F_0(k-m\omega^2)/D$, $B=F_0b\omega/D$, $D=(k-m\omega^2)^2+b^2\omega^2$. 无阻尼且 $\omega=\omega_0$ 时这组有限稳态不存在, 代入可得共振解含 $t\sin\omega_0t$, 振幅随时间增长.

扭摆的恢复矩为 $-\kappa\theta$, $\kappa$ 单位 N m/rad. 运动方程 $I\ddot\theta+\kappa\theta=0$, 周期 $2\pi\sqrt{I/\kappa}$, 不由地球重力决定. 复合物理摆可把各构件关于同一轴的惯量相加, 恢复矩系数为 $g\sum_jM_jd_j$, 其中 $d_j$ 是平衡时向下的有向距离. 例如端悬均匀杆 $M,L$ 加末端点质量 $m$, $I=ML^2/3+mL^2$, 重力矩系数 $g(ML/2+mL)$, 所以周期 $2\pi\sqrt{(ML^2/3+mL^2)/[gL(M/2+m)]}$. 若构件在支点上方, 对恢复系数贡献为负, 须核对总系数为正.

\originalsection{把不同机制化为振子}
圆柱质量 $M$、半径 $R$, 其轴心接水平弹簧 $k$, 在粗糙水平面上无滑小幅往复. 用轴心偏移 $x$ 描述, $\omega=\dot x/R$, 故总动能为 $(M+I/R^2)\dot x^2/2$, 势能为 $kx^2/2$. 能量微分给 $(M+I/R^2)\ddot x+kx=0$. 均匀实心柱 $I=MR^2/2$, 所以角频率 $\sqrt{2k/(3M)}$, 比同质量滑块低; 转动消耗的动能表现为更大有效质量. 此推导要求静摩擦足够, 若打滑须改用平动、转动两方程.

等截面 U 管中不可压缩液柱长度 $L$, 截面 $A$, 密度 $\rho$, 两竖直管口液面初等高. 推动后某侧升 $x$, 另一侧降 $x$, 整柱速率 $\dot x$. 相比平衡位置, 将体积 $Ax$ 提高平均高度 $x$ 得势能增加 $\rho Agx^2$; 动能为 $\rho AL\dot x^2/2$. 无黏性、等截面且两端同压时 $\rho AL\ddot x=-2\rho Agx$, 角频率 $\sqrt{2g/L}$. 这是原第 23 章以机械能处理的小振动例, 仅用质量和势能, 不扩写流体课程. 两个例子说明恢复作用未必来自一条单独弹簧, 而“振动质量”也可能包含转动或整段连续物质.

弱阻尼下还可从共振峰半高宽定义品质因子. 在 $|\omega-\omega_0|\ll\omega_0$ 展开 $k-m\omega^2\simeq-2m\omega_0(\omega-\omega_0)$, 平均吸收功率为以 $\omega_0$ 为中心的近似洛伦兹峰, 半高处 $|\omega-\omega_0|=b/(2m)$. 全半高宽 $\Delta\omega\simeq b/m$, 所以 $Q\simeq\omega_0/\Delta\omega=m\omega_0/b$. 该谱宽解释要求弱阻尼, 不能把任意阻尼的位移峰宽直接等同此式. 上述复数工具中 $z=x+iy$ 的共轭为 $\bar z=x-iy$, 模 $|z|=\sqrt{z\bar z}=\sqrt{x^2+y^2}$, 辐角按 $(x,y)$ 象限确定; 因此取驱动复幅值的模和相位有明确的几何意义.
